\section{Week - 08}
\begin{physicsmodule}{Feb 15 - Feb 21}{RoyalBlue}
  \begin{lawbox}
    \subsection*{\sffamily\bfseries Week - 08}
    \begin{enumerate}
    \item Complete the reading of the papers:
    \begin{enumerate} [\ding{212}]
        \item \textbf{Magnetic field-induced non-trivial electronic topology in} $\mathbf{Fe_{3-x}Ge Te_2}$ \href{https://doi.org/10.1063/5.0052952}{\textit{Appl. Phys. Rev. 8, 041401 (2021)}}
        \item \textbf{Magnetocrystalline Anisotropy Enables Field-Free Deterministic Switching in $\mathrm{Tm_3Fe_5O_{12}/Pt}$ Bilayers: An Atomistic Spin Dynamics Study} \href{https://doi.org/10.1021/acsami.5c23180}{\textit{ACS Applied Materials \& Interfaces}, 2026} - \textcolor{blue}{Revise if already read}
        \item Revisit the Exchange bias-related paper from Ahsan's paper and other VAMPIRE papers, and list most of the papers related to our work on exchange bias.
    \end{enumerate}
    \item Work on writing the text for the system description for the paper.
    \item Complete the thickness dependence text writing by Thursday night.
    \item Meet Emmanuel when he emails you and explain to him in detail about the exchange-bias related work for the short paper, and we will finalize what properties to vary and see the plots.
    \item 
\end{enumerate}
  \end{lawbox}
\end{physicsmodule}
%%%%%%%%%%%%%%%%%%%%%%%%%%%%%%%%%%%%%%%%%%%%%%%%%%%%%%%%%%%%%%%%
\subsection*{\underline{Sunday Feb 15, 2026}}
\subsection*{To do Tasks:}
\begin{enumerate}
    \item Complete the week plan and stick to it. 
    \item Plan at least 3--notes, which I will be doing before the end of February. 
    \begin{itemize}
        \item Zeroth Law of Thermodynamics and First Law of Thermodynamics together.
        \item Coulomb's Law and some basics of that portion. [Go and see this on Tuesday in the  department.]
    \end{itemize}  
\end{enumerate}
\subsection*{Focused Tasks of the Day}
\begin{enumerate}[\ding{212}]  % You can change the number to any dingbat symbol you want
    \item 
\end{enumerate}
\subsection*{Notes:}
\begin{enumerate}
    \item 
\end{enumerate}
%%%%%%%%%%%%%%%%%%%%%%%%%%%%%%%%%%%%%%%%%%%%%%%%%%%%%%%%%%%%%%%%
\subsection*{\underline{Monday Feb 16, 2026}}
\subsection*{To do Tasks:}
\begin{enumerate}
    \item Go to the department in the morning and complete the printing and pasting of the GSA flyers throughout the APSC premises. \cmark
    \item Complete the reading of the exchange-bias papers and make the important points to consider and discuss. 
    \item Make different folders of paper to share with Emmanuel, and ask him to add the papers to the same folder if he finds any relevance to our project.
    \item Add the references provided by Rajnandini to her new Magnetism note. \cmark 
\end{enumerate}
\subsection*{Focused Tasks of the Day}
\begin{enumerate}[\ding{212}]  % You can change the number to any dingbat symbol you want
    \item Complete the laundry thing and go to the department in the afternoon.
    \item Read the exchange-bias papers.
\end{enumerate}
\subsection*{Notes:}
\begin{enumerate}
    \item Read July 22, 2025, Google note page carefully for parameters and system details to be used in the paper.
    \item Make the change in the interface damping and see the hysteresis plot, if possible, see the change in the M-H loop area, and if possible, $H_c$ also.
\end{enumerate}
%%%%%%%%%%%%%%%%%%%%%%%%%%%%%%%%%%%%%%%%%%%%%%%%%%%%%%%%%%%%%%%%%
\subsection*{\underline{Tuesday Feb 17, 2026}}
\subsection*{To do Tasks:}
\begin{enumerate}
    \item Complete the reading of the exchange-bias papers and make the important points to consider and discuss. 
    \item Make different folders of paper to share with Emmanuel, and ask him to add the papers to the same folder if he finds any relevance to our project.
    \item Start the writing of thickness dependent part.
    \item Make the storyline up and other main points details for the GRS presentation.
\end{enumerate}
\subsection*{Focused Tasks of the Day}
\begin{enumerate}[\ding{52}]  % You can change the number to any dingbat symbol you want
    \item Complete the first Exchange bias-related paper of Ahsan, EB\_02.
\end{enumerate}
\subsection*{Notes:}
\begin{enumerate}
    \item Spent a good amount of time creating the page theme; the old one is taking too long. 
    \item I have to ask some of my old contacts for a new website theme.
\end{enumerate}
% %%%%%%%%%%%%%%%%%%%%%%%%%%%%%%%%%%%%%%%%%%%%%%%%%%%%%%%%%%%%%%%%%
\subsection*{\underline{Wednesday Feb 18, 2026} }
\subsection*{To do Tasks:}
\begin{enumerate}
    \item Complete the note-taking of the second exchange bias paper of Junaid. \cmark
    \item List the main and initial points that are to be done by Emmanuel, and ask him to be super active in the work. 
    \item Let him understand the Hysteresis loop first, and its applications, and why that is fundamental for all of the magnetic measurements.
    \item Details things to be said to Emmanuel for hysteresis calculations:
    \begin{itemize}
        \item 
    \end{itemize}
\end{enumerate}
\subsection*{Focused Tasks of the Day}
\begin{enumerate}[\ding{52}]  % You can change the number to any dingbat symbol you want
    \item Complete the note taking for the $2^{\mathrm{nd}}$ EB paper of Junaid.\cmark
    \item Start revising the thickness-dependent part of the paper.
    \item Think and work on the storyline for the GRS talk.
\end{enumerate}
\subsection*{Notes:}
\begin{enumerate}
    \item There is this paper of SS Parkin for $\mathrm{Co_{0.9}Fe_{0.1}}/IrMn$ where there is a plot in which $H_{EB}$ vs $t_{FM}$ is plotted, and it is useful to our work.
    \item The Paper [\href{https://doi.org/10.1103/PhysRevB.111.184416}{\textit{Phys. Rev. B} \textbf{111}, 184416 (2025)}] talks a lot about the range of the interface exchange coupling, and the range taking is really helpful in our work for Project 2.
    \item Skim this paper, \href{https://doi.org/10.1002/adma.202311643}{Adv. Mater. 2024, 2311643}: For CoGd/IrMn (Atomistic done in this work.)
    \item See the paper: \href{https://doi.org/10.1016/j.jmmm.2023.170680}{Analysis of ultrafast magnetization switching dynamics in exchange-coupled ferromagnet-ferrimagnet heterostructures.}
    \item 
\end{enumerate}
% %%%%%%%%%%%%%%%%%%%%%%%%%%%%%%%%%%%%%%%%%%%%%%%%%%%%%%%%%%%%%%%%%
\subsection*{\underline{Thursday Feb 19, 2026}}
\subsection*{To do Tasks:}
\begin{enumerate}
    \item See some of the videos of the PPT animation and see if you can find a gif to depict the precession dynamics. [\textcolor{blue}{Partially done}]
    \item Collect the images that you want to show in the presentation. \xmark
    \item Make a separate note for each of the page if you managed to get the diagrams. \cmark
\end{enumerate}
\subsection*{Focused Tasks of the Day}
\begin{enumerate}[\ding{52}]  % You can change the number to any dingbat symbol you want
    \item Watch the PPT animation things. 
    \item Work on the storyline for the talk, the talk time is 12 min + 3 min scheduled at 4:15 pm on Thursday. 
\end{enumerate}
\subsection*{Notes:}
\begin{enumerate}
    \item I was able to make an animation of LLG dynamics using Python and the matplotlib library, and got the code from the Gemini.
\end{enumerate}
% %%%%%%%%%%%%%%%%%%%%%%%%%%%%%%%%%%%%%%%%%%%%%%%%%%%%%%%%%%%%%%%%%
\subsection*{\underline{Friday Feb 20, 2026}}
\subsection*{To do Tasks:}
\begin{enumerate}
    \item Work on the re-reading of thickness-related information and note the details to be said at the GRS-2026 presentation and get the plots for the presentation. \xmark
    \item Start working on arranging the plots for the presentation. [\textcolor{blue}{Partially done}]
    \item Go for a short run, around 5 km, when the weather settles. \xmark
\end{enumerate}
\subsection*{Focused Tasks of the Day}
\begin{enumerate}[\ding{52}]  % You can change the number to any dingbat symbol you want
    \item Re-read the thickness-dependent reasoning and save the main points.
    \item Collect the images of the data and put them in a demo presentation. \cmark
\end{enumerate}
\subsection*{Notes:}
\begin{enumerate}
    \item 
\end{enumerate}
%%%%%%%%%%%%%%%%%%%%%%%%%%%%%%%%%%%%%%%%%%%%%%%%%%%%%%%%%%%%%%%%%
\subsection*{\underline{Saturday Feb 21, 2026} }
\subsection*{To do Tasks:}
\begin{enumerate}
    \item Go for a run and work on the easy run, and do not hurt your shins and calves. Went to run on Jamestown Island. \cmark
    \item Complete the PPT for the presentation by the end of the day.
    \item 
\end{enumerate}
\subsection*{Focused Tasks of the Day}
\begin{enumerate}[\ding{212}]  % You can change the number to any dingbat symbol you want
    \item Make the storyline for the complete presentation. 
\end{enumerate}
\subsection*{Notes:}
\begin{enumerate}
    \item 
\end{enumerate}
\stardivider
%%%%%%%%%%%%%%%%%%%%%%%%%%%%%%%%%%%%%%%%%%%%%%%%%%%%%%%%%%%%%%%%%
